% Final section; all metrics and model-complexity values were checked against
% the canonical comparison and model-specific experiment artifacts.

\section{Comparison of Results}
\label{sec:g}

\subsection{Overall comparison table}

\begin{table}[H]
  \centering
  \caption{All 5 models on the shared 885-sample held-out test set}
  \label{tab:g1}
  \resizebox{\textwidth}{!}{%
  \begin{tabular}{@{}lrrrrrrrrr@{}}
    \toprule
    Model & Test acc & Error & F1 macro & ROC-AUC & Rec.\ Dropout & Rec.\ Enrolled & Rec.\ Graduate & Depth & Leaves \\
    \midrule
    M0 (baseline)        & 0.6689 & 0.3311 & 0.6083 & 0.7195 & 0.6796 & 0.3836 & 0.7647 & 27 & 634 \\
    M1 (pruning)          & \textbf{0.7559} & \textbf{0.2441} & \textbf{0.6729} & \textbf{0.8477} & 0.6866 & 0.3459 & \textbf{0.9480} & \textbf{5} & \textbf{17} \\
    M2a (class\_weight)   & 0.6508 & 0.3492 & 0.5854 & 0.7053 & 0.6796 & 0.3396 & 0.7443 & 28 & 696 \\
    M2b (SMOTE)           & 0.6881 & 0.3119 & 0.6387 & 0.7421 & \textbf{0.7077} & \textbf{0.4654} & 0.7557 & 39 & 847 \\
    M3 (early warning)    & 0.5412 & 0.4588 & 0.4931 & 0.6254 & 0.5317 & 0.3270 & 0.6244 & 30 & 963 \\
    \bottomrule
  \end{tabular}%
  }
\end{table}

\begin{figure}[H]
  \centering
  \includegraphics[width=\textwidth]{comparison.png}
  \caption{Four-panel comparison of M0-M1-M2a-M2b-M3 on the same
    held-out test set (885 samples).}
  \label{fig:g1}
\end{figure}

Reading the four panels: top-left (test accuracy and error rate) shows
M1 highest and M3 lowest; top-right (the four macro metrics: Precision,
Recall, F1, ROC-AUC) shows M1 leading on all four; bottom-left (per-class
recall) shows M2b dominating on Dropout and Enrolled while M1 dominates
on Graduate; bottom-right (tree complexity, leaf count on a log scale)
shows M1 far more compact than the other four models.

\subsection{Comparing each improvement against the baseline}

\textbf{M1 (Pruning) improves almost every aggregate metric.} Test
accuracy rises from 0.6689 to 0.7559 (\textbf{+8.70 percentage points}),
error rate drops from 33.11\% to 24.41\%, F1 macro rises 6.47 points, and
ROC-AUC macro rises 12.82 points. At the same time, complexity drops
sharply: depth 27$\to$5, leaves 634$\to$17 (\textbf{about 37$\times$
fewer}). This is a favorable result: aggregate performance and structural
simplicity both improve, although Enrolled recall decreases slightly
(0.3836$\to$0.3459, $-3.77$ points). Accuracy was the prespecified CV
criterion for selecting \texttt{ccp\_alpha}/\texttt{max\_depth}/
\texttt{min\_samples\_leaf}; because that criterion does not explicitly
protect minority-class recall, it can favor the majority Graduate class
whenever doing so raises the total number of correct predictions.

M2 (Class Imbalance) gives two opposite results, consistent with the
experimental nature of this improvement. M2a
(\texttt{class\_weight='balanced'}) \textbf{does not improve} M0 on this
test set: test accuracy drops 1.81 points, and even Enrolled recall, the
metric this technique targets, also drops (0.3836$\to$0.3396). In
contrast, M2b (SMOTE on train only) improves both overall accuracy
(+1.92 points) and the recall of the two hardest classes: Dropout recall
+2.82 points, Enrolled recall \textbf{+8.18 percentage points} (61/159
$\to$ 74/159 Enrolled students correctly identified), at the cost of a
slight 0.90-point drop in Graduate recall. M2b also raises Enrolled
precision relative to M0 (0.3567$\to$0.4044), so the higher recall does
not simply come from the model predicting Enrolled more often. Vanilla
SMOTE on encoded categorical features still has the quantitative limitations
documented in Section~\ref{sec:f2-limitation}, so this result is best read as
an improvement on this particular held-out set rather than a guarantee for a
new cohort or institution.

\textbf{M3 (Early Warning) drops performance substantially, and this is
the expected result, not a failure.} Test accuracy falls from 0.6689 to
0.5412 (\textbf{$-12.77$ percentage points}), F1 macro falls 11.52
points. A model-grounded explanation comes from M0's own tree rather than
the original paper: by M0's Gini feature importance, 3 of the top
5 features are \texttt{2nd sem approved} (rank~1, importance 0.340, over
a third of the whole tree's total importance), \texttt{2nd sem grade}
(rank~2, 0.045), and \texttt{2nd sem enrolled} (rank~4, 0.045). All
three fall inside the 12 columns dropped in M3 because they are only
known \emph{after} the second semester ends. The other two features in
M0's real top-5, \texttt{Tuition fees up to date} (rank~3) and
\texttt{Admission grade} (rank~5), are retained in M3 because they are
treated as available at enrollment under this study's feature-availability
assumption. The importance ranking is associative rather than causal, but the
held-out performance gap is consistent with the loss of
\texttt{2nd sem approved}, which alone accounts for more importance than
the next four features combined. Under that assumption, M3 uses only
variables treated as available at enrollment (program of study, admission
grade, tuition status, the 3 macroeconomic variables, \ldots); production use
would require a data dictionary and snapshot timestamps to verify that timing.

\begin{quote}
\small
\emph{The original paper ranks first-semester approvals second using
permutation importance from four ensemble algorithms
\citep{realinho2022predicting}. That is a different method from this
project's M0 Gini importance. The feature ranks ninth in M0 (importance
0.029), so the different rankings are expected. The paragraph above uses
only importance values measured from M0 itself.}
\end{quote}

\subsection{Which model is ``best''?}

There is no single answer: it depends on the criterion.

\begin{table}[H]
  \centering
  \caption{Best model by criterion}
  \label{tab:g2}
  \begin{tabular}{@{}p{4.3cm}lp{6.3cm}@{}}
    \toprule
    Criterion & Best model & Why \\
    \midrule
    Overall accuracy / F1 / ROC-AUC & \textbf{M1} & Highest on nearly every aggregate metric, and the most compact tree \\
    Detecting Dropout students & \textbf{M2b} & Highest Dropout recall (0.7077) \\
    Detecting Enrolled students (hardest class) & \textbf{M2b} & Highest Enrolled recall (0.4654), $\sim$1.21$\times$ M0 \\
    Interpretability / presentation & \textbf{M1} & Only 17 leaves, fully readable, vs.\ 634--963 for the other models \\
    Early deployment under the stated timing assumption & \textbf{M3} & The only model that excludes all semester-1/2 outcomes; feasibility remains conditional on verifying retained-field timestamps \\
    \bottomrule
  \end{tabular}
\end{table}

\textbf{On accuracy alone, M1 wins.} But the main lesson of comparing all
five model configurations is that \textbf{``best'' depends on the application goal.} M1
is best for an accurate, interpretable classification system used once
semester data already exists. M2b is best if the goal is to miss as few
at-risk students as possible (Dropout/Enrolled), accepting lower overall
accuracy than M1 (0.6881 vs.\ 0.7559) in exchange for the better Dropout
and Enrolled recall shown above; M2b still improves on M0's accuracy, so
this trade-off is only relative to M1, not to the unimproved baseline.
Under this study's feature-availability assumption, M3 is the only candidate
for \textbf{intervention at enrollment} because it excludes every semester
outcome. The other four configurations require semester data; M3's own
feasibility remains conditional on verifying when each retained field is
actually captured.
